% ============================================================
% 2 相关工作
% ============================================================
\section{相关工作}

\subsection{深度伪造检测}

深度伪造检测方法大致可分为基于手工伪影（blending artifacts、频域指纹、生理信号）
与基于深度特征两大类。FaceForensics++\cite{rossler2019faceforensics} 建立了大规模基准数据集
与评测范式，此后大量工作围绕更强的判别网络展开。这类方法回答的是“真或假”，
输出为二值标签；在取证场景中，仅给出真假标签往往不足以定位内容来源，
这也正是主动溯源工作的出发点。

\subsection{主动检索与哈希溯源}

HashVAE\cite{liu2026hashvae} 是图像级的代表工作：先把待查图分为前景与背景，
再由多尺度多分支残差主干分别编码为哈希码，前景码额外接检测模块；
训练采用无监督重构、最小熵比特约束与身份感知四元组损失，推理时以两步检索缩小候选池。
在 FaceForensics++ 上，其报告的最高设定（$1024$ bit、c23）Top-5 为 $98.23\%$。

ViTHash（HashViT）\cite{pei2021vithash} 把溯源做到视频级：以视觉 Transformer 学习哈希中心，
并设计 Hash Triplet Loss 处理“同一场景不同角度、相似场景同一人物”的极难负样本；
检索到源视频后，再由定位模块对篡改区域做解释。该工作同时覆盖换脸、视频修复与视频拼接。

两项工作都把溯源表述为哈希检索，并在各自设定的闭集、干净、每视频单帧条件下取得很高读数。
它们为本文提供了任务范式与参照读数。

\subsection{零训练描述子与近重复检索}

感知哈希（pHash）、缩略图与颜色直方图是近重复检索中的经典零训练描述子；
冻结的通用视觉主干（如 DINOv2\cite{oquab2024dinov2}）则可提供更强的语义特征。
在“同源背景”这一任务上，这些描述子无须任何训练即可工作，
因而构成衡量学习方法真实收益的自然参照。本文把它们纳入统一协议，
并据此确定帧级强度基线（第 5.2 节）。

\subsection{几何验证}

经典图像配准以 SIFT\cite{lowe2004sift} 特征匹配配合 RANSAC\cite{fischler1981ransac}
做几何验证：前者提供尺度不变的关键点，后者以随机采样剔除误匹配。
这条路线精度高，但依赖特征检测器与随机采样流程。近几年也出现了以可微方式
估计几何变换与统计内点的模块，可直接嵌入网络。
本文在检索期尝试了后一条路径：以软匹配代替硬匹配、以可微 DLT 代替随机采样、
以软内点占比代替硬内点计数，并与 SIFT$+$RANSAC 在同一几何退化协议下并列比较（第 4.3、5.6 节）。

\subsection{本文的位置}

在上述脉络之上，本文把评测条件补齐到退化与开集两类真实场景，
在同一协议下同时报告零训练描述子、学习编码器与两种几何验证路线，
使“学习带来多少收益”“哪种几何验证更合适”这两个问题可以在同一张表上回答。
